\section*{Sets and Indices}

\begin{itemize}
    \item Years in the planning horizon: $T = \{1,2\}$.
\end{itemize}

\section*{Parameters}

\begin{itemize}
    \item $a_1$ (code: \texttt{a1}): fighter jets produced in year 1.
    \item $a_2$ (code: \texttt{a2}): fighter jets produced in year 2.
    \item $c$ (code: \texttt{training\_capacity}): number of pilots that one training jet can train in one year.
    \item $D$ (code: \texttt{training\_duration}): duration of the training program in years. In this instance, $D=2$.
\end{itemize}

For Instance 2, the data are
\[
a_1=10,\qquad a_2=15,\qquad c=5,\qquad D=2.
\]

\section*{Decision Variables}

The implemented logic comments describe the following variables:

\begin{itemize}
    \item $x_1$: number of jets allocated to training in year 1.
    \item $x_2$: number of training jets available in year 2 (including year-1 training jets that continue training).
\end{itemize}

However, the current implementation does not optimize over these variables; instead it directly sets the reported quantities using parameter values only.

\section*{Objective}

The search logic effectively evaluates
\[
\max \; c\,(a_1+a_2),
\]
which for this instance gives
\[
\max \; 5(10+15)=125.
\]

Equivalently, in terms of the commented variable interpretation, the intended objective expression is
\[
\max \; c(x_1+x_2),
\]
but the code substitutes the specific choice
\[
x_1=a_1,\qquad x_2=a_2,
\]
and therefore computes
\[
c(x_1+x_2)=c(a_1+a_2).
\]

\section*{Constraints}

The source code contains only commented model logic and no enforced optimization constraints. The commented logic states:

\begin{align}
x_1 &\le a_1,\\
x_2 &\le a_1+a_2,\\
x_1 &\le x_2.
\end{align}

In the actual implementation, these constraints are not instantiated or solved. Instead, the code directly reports
\[
x_1=a_1,\qquad x_2=a_2,
\]
so the computed total pilots is
\[
c(a_1+a_2).
\]

\section*{Notes on Implementation}

\begin{itemize}
    \item The file \texttt{current\_heuristic.py} does not build an LP, MIP, or any other optimization model.
    \item No solver is called. The objective value is computed explicitly as
    \[
    \texttt{total\_pilots}=\texttt{training\_capacity}\times(\texttt{a1}+\texttt{a2}).
    \]
    \item The printed lines label year-1 training jets as \texttt{a1} and year-2 training jets as \texttt{a2}. Thus the implementation treats all year-1 production and all year-2 production as contributing once to pilot training, without enforcing jet availability linkage beyond this direct arithmetic.
    \item Although the comments describe a 2-year training-allocation problem with variables $(x_1,x_2)$, the executable logic is simply a closed-form evaluation:
    \[
    \text{Objective value} = c(a_1+a_2)=125.
    \]
\end{itemize}
